\documentclass[11pt,a4paper]{article}

\usepackage[T1]{fontenc}
\usepackage[utf8]{inputenc}
\usepackage{lmodern}
\usepackage[margin=2.3cm]{geometry}
\usepackage{amsmath,amssymb}
\usepackage{booktabs}
\usepackage{makecell}
\usepackage{graphicx}
\usepackage{pdflscape}
\usepackage{caption}
\usepackage{enumitem}
\usepackage{microtype}
\usepackage[hidelinks]{hyperref}

\captionsetup{font=small,labelfont=bf}
\renewcommand{\theadfont}{\footnotesize\bfseries}
\setlist{itemsep=0.35em}
% report tables are lettered A--D so they cannot be confused with the paper's Table~1
\renewcommand{\thetable}{\Alph{table}}

% notation
\newcommand{\Mnd}{M_{\mathrm{nd}}}
\newcommand{\kstag}{k_{\mathrm{stag}}}
\newcommand{\dx}{\Delta x}
\newcommand{\maxdx}{\max|\Delta x_e|}
\newcommand{\relL}{\|\Delta x\|_2/\|x\|_2}
\newcommand{\RMS}{\mathrm{RMS}}
\newcommand{\me}[2]{$#1\times#2$}
\newcommand{\file}[1]{\texttt{#1}}

\title{One stopping criterion for the three methods:\\[0.3em]
  \large why per-iteration design-change norms cannot be made universal,\\
  and what to use instead}
\author{}
\date{9 October 2026}

\begin{document}
\maketitle

\paragraph{Question.}
The nine-mesh Table~1 campaign uses a different stopping rule per method
($\maxdx$ for Yuksel and the proposed approach,
$\|\dx\|_2 < 0.05\sqrt{n_e/3200}$ for Du--Olhoff). A common rule is needed, but
every candidate tried so far ($\maxdx$, $\relL$) has the same defect: a tolerance
that lets \me{800}{100} converge cleanly produces long flat tails on the coarse
meshes, and a tolerance that stops \me{240}{30} at ${\sim}50$ iterations leaves
\me{800}{100} gray.

\paragraph{Answer in one paragraph.}
This is not a tolerance-selection problem; it is a metric problem. All
per-iteration design-change norms (max, $L_2$, RMS, relative $L_2$) measure how
fast the optimizer is still \emph{moving elements}, not how far the design is
from its final state. Once $\omega_1$ and the discreteness $\Mnd$ have settled,
all three optimizers keep moving boundary elements by 0.01--0.10 per iteration
(OC and the adaptive-box MMA both oscillate at the boundary), so the norms sit
on a \emph{noise floor} that is 3--10$\times$ above the thresholds in use, and
then decay slowly and non-monotonically (Table~\ref{tab:B},
Figs.~\ref{fig:traces-proposed}--\ref{fig:traces-olhoff}). A stop therefore
happens when the noise happens to dip below $\varepsilon$: with $\varepsilon$
near the floor the dip comes early and may be a mid-run plateau (the gray
\me{800}{100} Olhoff design at $c = 0.2$ stopped at outer iteration~37); with
$\varepsilon$ below the floor the dip comes hundreds of iterations after
stagnation (\me{240}{30} Proposed at $\relL < 10^{-3}$: iteration~271 versus
stagnation at ${\sim}45$). Both failure modes are properties of the metric, so
no $\varepsilon$ fixes both, and the ``universal'' $\varepsilon$ does not exist
for any of these norms. The remedy is to stop on \emph{stagnation of the
quantities the paper actually reports}, measured over a window:
\textbf{stop at the first iteration $k \ge W$ at which, over the last $W = 10$
iterations, the method's own objective varied by less than 0.1\,\% and $\Mnd$
varied by less than 0.005} (rule R1 below). R1 is dimensionless, contains no
$n_e$, is identical for all three methods, and in an offline replay on 23
recorded histories (three methods, \me{160}{20} \ldots\ \me{800}{100}) it stops
within 0.48--2.68$\times$ of the measured stagnation iteration with an $\omega_1$
deficit of at most 0.91\,\% relative to running 300 (Proposed, Olhoff) or 600
more (Yuksel stage~2) iterations. A windowed $\maxdx$ rule (R2) is a defensible
second choice, but it is not universal: it never fires for the proposed method
at \me{800}{100} because single OC elements keep jumping by more than 0.04.

\medskip\noindent
No code was changed for this study. Everything below was produced with the
existing solvers and drivers through recording runs placed in a scratch
directory; scripts, per-iteration histories (CSV) and figures are in
\file{data/\allowbreak{}}, \file{fig/\allowbreak{}} and \file{scripts/\allowbreak{}} next to this report.

%======================================================================
\section{What was measured}

\paragraph{Recording runs.}
Each method was run on the simply supported beam of Table~1 with its
production configuration (frozen profiles
\file{proposed\_\allowbreak{}practical\_\allowbreak{}move02\_\allowbreak{}tol001},
\file{yuksel\_\allowbreak{}practical\_\allowbreak{}move01\_\allowbreak{}tol001}, production Olhoff preset
\file{duOlhoffPedersenAdaptiveBoxSensitivityFiltered}) but with the native stop
\emph{disabled or extended}, so that the whole trajectory past the usual stop is
on record (table below).

\begin{table}[htb]
\centering\small
\caption*{Recording runs.}
\begin{tabular}{@{}l p{3.0cm} p{4.6cm} p{5.0cm}@{}}
\toprule
method & meshes & horizon & recorded per iteration \\
\midrule
Proposed & \me{160}{20} \ldots\ \me{800}{100} (9)
  & 300 SIMP iterations (production stop at $\maxdx \le 0.01$ recorded but not acted on)
  & $x_{\mathrm{phys}}$; $\maxdx$, $\relL$, $\RMS(\dx)$, $\Mnd$, objective,
    E1 structural $\omega_1$ (every iteration up to \me{320}{40}, every 5th from \me{400}{50}) \\
\addlinespace
Yuksel & \me{160}{20} \ldots\ \me{800}{100} (9)
  & stage~1 on its native rule ($\maxdx < 0.01$, budget 2000); stage~2 extended to 600 iterations
  & same, E1 $\omega_1$ every 5th iteration \\
\addlinespace
Du--Olhoff & \me{160}{20}, \me{240}{30}, \me{320}{40}, \me{400}{50}, \me{800}{100}
  & 300 outer iterations, stop disabled
  & $\rho$ replayed from the recorded increments; same metrics, native $\omega_1$ \\
\bottomrule
\end{tabular}
\end{table}

The stopping rule of every method is a pure \emph{observer} of the trajectory
(it changes when the loop exits, not what the loop does), with one exception:
Yuksel's stage-1 tolerance decides the design handed to stage~2. Hence every
candidate rule can be replayed exactly on these records. This was checked: the
replay reproduces the production stop iterations 107/236/207/182 for Proposed
and 121/111/101/93 for Olhoff at $c = 0.05$, and 58/52/67/62 at $c = 0.2$ --- the
numbers of the \file{campaign\_\allowbreak{}mac\_\allowbreak{}convergence\_\allowbreak{}*} campaigns.

\paragraph{Ground truth for ``the design has stopped changing''.}
For each record, $\kstag$ is the first iteration after which $\omega_1$ (E1
structural $\omega_1$ for Proposed/Yuksel, native for Olhoff) stays within
0.5\,\% and $\Mnd$ within 0.01 of the end of the record. A rule is judged by
where it stops relative to $\kstag$ and by the $\omega_1$ it leaves on the table
relative to the end state.

\paragraph{Caveats of the record.}
\begin{enumerate}[label=(\roman*)]
\item Yuksel's stage~2 is still improving slowly at the 600-iteration cap on
  the finer meshes, so its ``end state'' is a cap, not a converged state; the
  losses quoted for Yuksel are relative to that cap.
\item For Proposed and Yuksel the late trajectory \emph{lowers} $\omega_1$
  (Table~\ref{tab:B}: Proposed \me{800}{100} peaks at 163.4\,rad/s and ends at
  161.4; Yuksel \me{160}{20} peaks at 158.3 and ends at 156.9), because their
  objective is a static surrogate, not $\omega_1$; negative ``losses'' in the
  tables are this effect.
\item One load case, one starting design, one filter radius per method;
  thresholds were chosen after seeing these data, so they are calibrated on
  this benchmark and should be preregistered, not re-tuned, for any other case.
\end{enumerate}

%======================================================================
\section{Diagnosis: why no \texorpdfstring{$\varepsilon$}{epsilon} works for a per-iteration change norm}

Figures~\ref{fig:traces-proposed}--\ref{fig:traces-olhoff} show the four
quantities per iteration for every mesh; Table~\ref{tab:B} gives their levels at
stagnation and at the end of the record.

\begin{enumerate}
\item \textbf{The norms do not go to zero when the design stops changing.}
  At $\kstag$ the median $\maxdx$ is 0.014--0.074 (Proposed), 0.03--0.10
  (Yuksel) and 0.05--0.07 (Olhoff); the median $\relL$ is
  $2.5$--$7.5\cdot10^{-3}$, $0.9$--$4.8\cdot10^{-3}$ and
  $5.3$--$6.4\cdot10^{-3}$. These are the optimizer's boundary oscillation (OC
  with a move limit flips boundary elements back and forth; the Olhoff adaptive
  box does the same at the 0.002--0.1 scale --- the period-2 behaviour already
  documented for \me{160}{20}). They carry no information about $\omega_1$ or
  $\Mnd$, both of which are flat.

\item \textbf{A threshold below the floor is reached by chance.}
  With $\maxdx \le 0.01$ the proposed method stops at 107, 236, 207, 182, 219,
  256, never, 297, never (\me{160}{20} \ldots\ \me{800}{100}) although $\kstag$
  is 29--95: the stop is the first iteration at which no element happened to
  jump. The waiting time for that event grows with the number of boundary
  elements, i.e.\ with the mesh. This is the ``flat tail'' seen on \me{240}{30}
  at $\relL < 10^{-3}$ (stop at 271). The earlier finding that 0.019 and 0.020
  give 98 and 47 iterations at \me{320}{40} is the same lottery.

\item \textbf{A threshold near the floor catches mid-run plateaus.}
  Du--Olhoff's $\RMS(\dx)$ dips during the adaptive-box transitions well before
  the design is finished; at $c = 0.2$ ($\RMS < 3.5\cdot10^{-3}$) the campaign
  stopped \me{640}{80}/\me{720}{90}/\me{800}{100} at outer 38/33/37 with
  $\omega_1 = 156/154/153$\,rad/s, 7\,\% below the converged 165. The
  \me{800}{100} record (Fig.~\ref{fig:traces-olhoff}) shows what that run was
  interrupted in: Du--Olhoff on this mesh de-grays in two slow phases,
  $\Mnd = 0.53 \to 0.30$ (iterations 25\,$\to$\,100) and $0.26 \to 0.16$
  (175\,$\to$\,240), with $\omega_1$ flat at 163.6--163.8 between iterations 125
  and 175 while the second phase is being prepared; $\kstag = 232$, against
  34--59 on \me{160}{20}--\me{400}{50}. The gray design at ``iteration 125'' in
  the current figure set ($\omega_1 = 163.6$, $\Mnd = 0.275$) is therefore an
  unfinished run, not a mesh artefact, and every rule that fired on that
  plateau --- single-iteration RMS/relative-$L_2$ at $5\cdot10^{-3}$ (37),
  $\maxdx \le 0.04$ with persistence (105), the objective-only window (122) ---
  produced it. The rules that waited for the second phase (relative
  $L_2 < 10^{-3}$: 256; $c = 0.05$: 246; R1: 241) are the same rules that
  produce the 100-iteration flat tails on \me{160}{20}--\me{240}{30}, which is
  the dilemma of the opening question restated: the dip statistics are
  different on the two meshes, the physics is not.

\item \textbf{Scaling the tolerance with $n_e$ cannot repair this.}
  $\|\dx\|_2 < 0.05\sqrt{n_e/3200}$ is literally
  $\RMS(\dx) < 8.84\cdot10^{-4}$, a mesh-intensive quantity (so the ``tailored
  to 3200'' objection is cosmetic: write it as RMS). The relative norm
  $\relL \approx \RMS(\dx)/\sqrt{\overline{x^2}}$ is the same quantity up to a
  factor ${\approx}\,1.4$--$1.5$. Both inherit the floor and the lottery.
  Table~\ref{tab:A} shows the RMS rule stopping Proposed at 209--289 and the
  relative-$L_2$ rule at 219--294, for designs that stagnated at 30--95.

\item \textbf{Mesh dependence is therefore not a scaling law to be fitted; it
  is the statistics of dips.} Across meshes the floor itself is nearly flat
  (Table~\ref{tab:B}), while the time to the first dip below a given
  $\varepsilon$ is erratic (Proposed at $\varepsilon = 0.01$: from 107 to
  never). No monotone $\varepsilon(n_e)$ maps onto that.
\end{enumerate}

\begin{figure}[p]
\centering
\includegraphics[width=\textwidth]{fig/traces_proposed.png}
\caption{Proposed method: per-iteration $\maxdx$ (solver-native), $\relL$,
  $\Mnd$ and $\omega_1$ for every recorded mesh.}
\label{fig:traces-proposed}
\vspace{1.5em}
\includegraphics[width=\textwidth]{fig/traces_yuksel.png}
\caption{Yuksel method: same quantities as Fig.~\ref{fig:traces-proposed}.}
\label{fig:traces-yuksel}
\end{figure}

\begin{figure}[tbp]
\centering
\includegraphics[width=\textwidth]{fig/traces_olhoff.png}
\caption{Du--Olhoff method: same quantities as Fig.~\ref{fig:traces-proposed}.
  Note the two de-graying phases at \me{800}{100}.}
\label{fig:traces-olhoff}
\end{figure}

%======================================================================
\section{Candidate common rules, replayed}

Table~\ref{tab:C} and Fig.~\ref{fig:summary} give, per method and mesh, the stop
iteration and the $\omega_1$ deficit of each candidate; Table~\ref{tab:D}
summarises over all records.

\paragraph{R1 (recommended) --- windowed stagnation of objective and discreteness.}
With $f_k$ the method's own objective (Du--Olhoff: $\omega_1$ or $\lambda_1$;
Yuksel and Proposed: their compliance-type objective) and
$\Mnd^{(k)} = 4\cdot\operatorname{mean}\bigl(x_{\mathrm{phys}}(1-x_{\mathrm{phys}})\bigr)$,
stop at the first $k \ge W$ with
\begin{equation}
\label{eq:R1}
\frac{\max\limits_{k-W\le j\le k} f_j - \min\limits_{k-W\le j\le k} f_j}{|f_k|} < 10^{-3}
\quad\text{and}\quad
\max_{k-W\le j\le k} \Mnd^{(j)} - \min_{k-W\le j\le k} \Mnd^{(j)} < 5\cdot10^{-3},
\qquad W = 10.
\end{equation}

Properties: dimensionless; no $n_e$, no move limit, no filter radius in it; both
quantities are already computed or trivially available in every loop (no extra
eigensolve); it is a statement about the outcome the paper reports (frequency
and black-and-whiteness), which is what a reviewer wants to be sure has
converged. The objective window alone stops too early on the Olhoff plateau at
\me{800}{100} (iteration 122, 1.1\,\% of $\omega_1$ and half of the de-graying
still to come); the $\Mnd$ window alone stops too early when the design is
discrete but $\omega_1$ is still moving (Yuksel); the conjunction is what makes
it universal. It is also the only candidate that resolves the dilemma as
stated: it stops Du--Olhoff \me{240}{30} at 57 (stagnation 34) and
\me{800}{100} at 241 (stagnation 232) with the same two numbers in the rule.

Replay: R1 fired on all 23 of 23 histories. Stop iteration relative to
$\kstag$: 0.48--2.68 (Proposed), 0.81--1.23 (Yuksel stage~2), 1.04--1.68
(Du--Olhoff); $\omega_1$ at the stop relative to the end state:
$-0.24 \ldots +0.33$\,\% (Proposed), $-0.46 \ldots +0.91$\,\% (Yuksel),
$-0.13 \ldots +0.27$\,\% (Du--Olhoff); the worst case is Yuksel \me{480}{60}
($+0.91$\,\%). $\Mnd$ at the stop exceeds its end value by at most 0.010.
Per-mesh numbers: Table~\ref{tab:C}.

\paragraph{R2 --- persistence version of the current rule.}
$\maxdx \le 0.04$ for $W = 10$ consecutive iterations. Same $\omega_1$ behaviour
as R1 on the coarse meshes, but
\begin{itemize}
\item it never fires for Proposed at \me{800}{100} within 300 iterations (single
  elements keep jumping by 0.05--0.1);
\item it stops Du--Olhoff \me{800}{100} at 105 with $\Mnd = 0.29$ (gray,
  1.3\,\% of $\omega_1$ missing);
\item it never fires for Yuksel at \me{640}{80}--\me{800}{100} because some
  element moves by the full move limit 0.1 at every stage-2 iteration
  (Table~\ref{tab:B}: median $\maxdx = 0.100$ at the end of the record);
\item its threshold is 20\,\% of the move limit of one method and 40\,\% of
  another's.
\end{itemize}
Acceptable as a documented fallback, not as the common rule.

\paragraph{Rules not recommended.}
Single-iteration $\maxdx \le 0.025$ (the 2026-10-07 retune) is tied for
$\omega_1$ but is the dip lottery with a larger $\varepsilon$; the two
objective-only windows are fine for Proposed and Yuksel but let Du--Olhoff stop
on its $\omega_1$ plateau at \me{240}{30} with 0.5\,\% left; element
flip-fraction and $\Mnd$-only rules miss up to 2.6\,\% of $\omega_1$.

\begin{figure}[tbp]
\centering
\includegraphics[width=0.92\textwidth]{fig/stop_rules_summary.png}
\caption{Stop iteration and $\omega_1$ deficit of the current and candidate rules
  per method and mesh, with $\kstag$.}
\label{fig:summary}
\end{figure}

%======================================================================
\section{What this changes in the paper and the campaign}

\begin{enumerate}
\item \textbf{One sentence replaces three.} ``All three implementations, and
  both stages of the Yuksel--Yilmaz method, terminate on the same rule: the
  first iteration at which, over the preceding ten iterations, the method's
  objective has varied by less than 0.1\,\% and the discreteness measure $\Mnd$
  by less than 0.005.'' Add that the rule is a stagnation test, not an
  optimality test, and that it was verified offline against 300-iteration
  extended runs ($\omega_1$ within $x$\,\% at every mesh; cite the
  supplementary table).

\item \textbf{Yuksel's stage-1 tolerance is part of the algorithm, not of the
  benchmark.} Changing it from the published 0.01 to 0.04 (the current driver)
  changes which design stage~2 starts from, hence the method. Either keep the
  native stage-1 rule and apply R1 to stage~2 only, or apply R1 to both stages
  and say so; the current mix is the one thing a reviewer can legitimately call
  unfair.

\item \textbf{Fix the inconsistency already in the text.} Section~4.1 states
  $\maxdx < 10^{-3}$ for Yuksel and the proposed method, the frozen profiles use
  0.01, and the current driver uses 0.04; the Table~1 figures in the text (345
  outer iterations, 1083 and 263 iterations at \me{400}{50}) belong to an older
  campaign. Whatever rule is adopted, the text, the manifest and the table must
  quote the same one.

\item \textbf{The $\sqrt{n_e/3200}$ form should go} even if the Olhoff native
  rule were kept for a sensitivity row: present it as
  $\RMS(\dx) < 8.8\cdot10^{-4}$.

\item \textbf{Regenerating Table~1 is cheap.} Because the rules are observers,
  the iteration counts and the designs at the new stop follow from the recorded
  histories without rerunning anything (discovery pass); only the clean timing
  needs the two-pass protocol already written in the benchmark plan (run to the
  discovered $k^*$ with history and diagnostics off). Expected effect on the
  counts: Du--Olhoff 57--78 instead of 93--122 outer iterations on
  \me{160}{20}--\me{400}{50} (its cost is dominated by the tail of nested MMA
  solves) and ${\approx}\,240$ at \me{800}{100}, where the method genuinely
  needs them; Proposed 44--98 instead of 107--297; Yuksel stage~2 unchanged in
  order of magnitude. The nine-mesh scaling exponent of Du--Olhoff will steepen
  accordingly: its iteration count is flat up to \me{400}{50} and then grows,
  and that is a property of the method on fine meshes, not of the stop rule.

\item \textbf{Say what the tail does.} For the two static-surrogate methods the
  extended tail lowers $\omega_1$ by up to 1.2\,\% while making the design
  slightly more discrete; for Du--Olhoff the tail adds ${\le}\,0.3$\,\%. This
  supports stopping at stagnation and is worth one sentence in the response to
  Reviewer~3's point~12 (gray regions / incomplete convergence): the gray that
  remains at the stop is the sensitivity-filter band, not an unfinished run
  (shown by the extended records: $\Mnd$ changes by $< 0.01$ after $\kstag$ on
  every mesh).
\end{enumerate}

%======================================================================
\section{Open points}

\begin{itemize}
\item Du--Olhoff at \me{480}{60}--\me{720}{90} was not recorded (each
  stop-disabled run costs 1--3\,h); the replay covers
  \me{160}{20}--\me{400}{50} and \me{800}{100} for that method.
\item The thresholds $10^{-3}$ / $5\cdot10^{-3}$ / $W = 10$ were chosen on these
  data. The sensitivity to $W$ (10 vs 20) and to the objective tolerance
  ($10^{-3}$ vs $5\cdot10^{-4}$) is in Table~\ref{tab:D}: tighter settings cost
  10--50\,\% more iterations for ${\le}\,0.1$\,\% of $\omega_1$. They should be
  frozen before the other examples (clamped beam, building) are rerun and
  reported as such.
\item Yuksel's stage~2 at the fine meshes is the one place where
  ``stagnation'' and ``converged'' differ by up to 0.9\,\% in $\omega_1$ within
  600 iterations; if that matters, the honest statement is that the method has
  a slow tail, not that the rule is wrong.
\end{itemize}

%======================================================================
\appendix
\section{Tables from the replay}

\begin{landscape}
\begin{table}[p]
\centering\small
\caption{Where the current rules stop, versus where the design has actually
  stagnated. $\kstag$ = first iteration after which $\omega_1$ stays within
  0.5\,\% and $\Mnd$ within 0.01 of the end of the extended record (300
  iterations for Proposed and Du--Olhoff, stage-2 cap 600 for Yuksel). Entries:
  stop iteration ($\omega_1$ loss vs end state, \%). \textit{none} = the rule
  never fired within the record.}
\label{tab:A}
\setlength{\tabcolsep}{4pt}
\begin{tabular}{@{}lrrcccccc@{}}
\toprule
mesh & record & $\kstag$
  & $\maxdx \le 0.01$
  & \makecell{$\RMS(\dx) < 8.8\cdot10^{-4}$\\$(= 0.05\sqrt{n_e/3200})$}
  & $\relL < 10^{-3}$
  & $\maxdx \le 0.04$
  & \makecell{$\RMS(\dx) < 3.5\cdot10^{-3}$\\$(c = 0.2)$}
  & $\relL < 5\cdot10^{-3}$ \\
\midrule
\multicolumn{9}{@{}l}{\textit{Proposed}} \\
$160\times20$ & $299$ & $29$ & 107 $(-0.12)$ & \textit{none} & \textit{none} & 32 $(-0.13)$ & 43 $(-0.24)$ & 45 $(-0.22)$ \\
$240\times30$ & $299$ & $95$ & 236 $(-0.03)$ & 266 $(-0.01)$ & 271 $(-0.00)$ & 36 $(+0.31)$ & 42 $(+0.33)$ & 43 $(+0.33)$ \\
$320\times40$ & $299$ & $30$ & 207 $(-0.05)$ & 209 $(-0.05)$ & 219 $(-0.04)$ & 35 $(-0.26)$ & 40 $(-0.17)$ & 41 $(-0.16)$ \\
$400\times50$ & $299$ & $34$ & 182 $(-0.01)$ & 221 $(+0.01)$ & 249 $(+0.01)$ & 32 $(-0.66)$ & 38 $(-0.42)$ & 38 $(-0.42)$ \\
$480\times60$ & $299$ & $30$ & 219 $(+0.03)$ & 257 $(+0.01)$ & 277 $(+0.00)$ & 50 $(+0.01)$ & 50 $(+0.01)$ & 51 $(+0.01)$ \\
$560\times70$ & $299$ & $39$ & 256 $(-0.02)$ & 254 $(-0.02)$ & 280 $(-0.01)$ & 57 $(-0.06)$ & 52 $(-0.09)$ & 52 $(-0.09)$ \\
$640\times80$ & $299$ & $39$ & \textit{none} & 289 $(+0.00)$ & \textit{none} & 47 $(-0.18)$ & 46 $(-0.18)$ & 47 $(-0.18)$ \\
$720\times90$ & $299$ & $39$ & 297 $(+0.00)$ & 234 $(+0.02)$ & 285 $(+0.01)$ & 45 $(-0.31)$ & 50 $(-0.19)$ & 50 $(-0.19)$ \\
$800\times100$ & $299$ & $44$ & \textit{none} & 263 $(+0.02)$ & 294 $(+0.01)$ & 56 $(-0.18)$ & 55 $(-0.18)$ & 55 $(-0.18)$ \\
\midrule
\multicolumn{9}{@{}l}{\textit{Yuksel (stage 2)}} \\
$160\times20$ & $720$ & $179$ & 244 $(-0.17)$ & 245 $(-0.16)$ & 251 $(-0.15)$ & 151 $(-0.77)$ & 149 $(-0.88)$ & 149 $(-0.88)$ \\
$240\times30$ & $767$ & $201$ & 320 $(-0.07)$ & 270 $(-0.09)$ & 312 $(-0.07)$ & 220 $(-0.40)$ & 199 $(-0.39)$ & 199 $(-0.39)$ \\
$320\times40$ & $851$ & $394$ & 572 $(-0.02)$ & 507 $(+0.03)$ & 531 $(+0.00)$ & 311 $(+1.05)$ & 282 $(+1.77)$ & 283 $(+1.77)$ \\
$400\times50$ & $914$ & $414$ & 732 $(+0.03)$ & 669 $(+0.06)$ & 682 $(+0.06)$ & 400 $(+0.56)$ & 351 $(+1.42)$ & 351 $(+1.42)$ \\
$480\times60$ & $1185$ & $884$ & \textit{none} & 781 $(+0.72)$ & 782 $(+0.72)$ & 781 $(+0.72)$ & 639 $(+1.58)$ & 639 $(+1.58)$ \\
$560\times70$ & $1432$ & $1119$ & \textit{none} & 1045 $(+0.70)$ & 1068 $(+0.63)$ & 1073 $(+0.62)$ & 882 $(+1.72)$ & 883 $(+1.72)$ \\
$640\times80$ & $2141$ & $1749$ & \textit{none} & 1853 $(+0.28)$ & 1854 $(+0.28)$ & 2006 $(+0.09)$ & 1583 $(+1.78)$ & 1583 $(+1.78)$ \\
$720\times90$ & $1761$ & $1374$ & \textit{none} & 1383 $(+0.48)$ & 1491 $(+0.31)$ & 1491 $(+0.31)$ & 1212 $(+1.51)$ & 1212 $(+1.51)$ \\
$800\times100$ & $1749$ & $1329$ & \textit{none} & 1424 $(+0.33)$ & 1509 $(+0.19)$ & 1604 $(+0.12)$ & 1203 $(+1.42)$ & 1203 $(+1.42)$ \\
\midrule
\multicolumn{9}{@{}l}{\textit{Du--Olhoff}} \\
$160\times20$ & $300$ & $53$ & 122 $(-0.01)$ & 121 $(-0.01)$ & 122 $(-0.01)$ & 67 $(+0.25)$ & 58 $(+0.32)$ & 58 $(+0.32)$ \\
$240\times30$ & $300$ & $34$ & 89 $(-0.04)$ & 111 $(-0.02)$ & 114 $(-0.03)$ & 61 $(-0.04)$ & 52 $(-0.07)$ & 52 $(-0.07)$ \\
$320\times40$ & $300$ & $59$ & 85 $(-0.10)$ & 101 $(-0.08)$ & 128 $(-0.06)$ & 67 $(-0.22)$ & 67 $(-0.22)$ & 67 $(-0.22)$ \\
$400\times50$ & $300$ & $52$ & 79 $(-0.13)$ & 93 $(-0.08)$ & 105 $(-0.06)$ & 60 $(-0.33)$ & 62 $(-0.31)$ & 62 $(-0.31)$ \\
$800\times100$ & $300$ & $232$ & 245 $(-0.04)$ & 246 $(-0.04)$ & 256 $(-0.04)$ & 18 $(+9.27)$ & 37 $(+7.35)$ & 37 $(+7.35)$ \\
\bottomrule
\end{tabular}
\end{table}

\begin{table}[p]
\centering\footnotesize
\setlength{\tabcolsep}{3.5pt}
\caption{Level of the per-iteration change metrics once the design has
  stagnated. Median over the 20 iterations after $\kstag$ (``@stag'') and over
  the last 20 recorded iterations (``end''). These are the values a threshold
  must exceed to stop at stagnation, and the floor it must stay above to stop
  at all. $\omega_1$ in rad/s.}
\label{tab:B}
\begin{tabular}{@{}lr ccc ccc ccc cc@{}}
\toprule
 & & \multicolumn{3}{c}{@stag} & \multicolumn{3}{c}{end}
   & \multicolumn{3}{c}{$\omega_1$} & \multicolumn{2}{c}{$\Mnd$} \\
\cmidrule(lr){3-5}\cmidrule(lr){6-8}\cmidrule(lr){9-11}\cmidrule(l){12-13}
mesh & $\kstag$ & $\maxdx$ & rel-$L_2$ & RMS & $\maxdx$ & rel-$L_2$ & RMS
  & @stag & max & end & @stag & end \\
\midrule
\multicolumn{13}{@{}l}{\textit{Proposed}} \\
$160\times20$ & $29$ & $0.021$ & $0.0061$ & $0.0040$ & $0.006$ & $0.0021$ & $0.0014$ & $153.6$ & $153.9$ & $153.5$ & $0.261$ & $0.252$ \\
$240\times30$ & $95$ & $0.014$ & $0.0025$ & $0.0017$ & $0.003$ & $0.0005$ & $0.0004$ & $157.1$ & $157.7$ & $157.6$ & $0.171$ & $0.161$ \\
$320\times40$ & $30$ & $0.025$ & $0.0050$ & $0.0034$ & $0.002$ & $0.0004$ & $0.0002$ & $159.2$ & $159.2$ & $158.7$ & $0.129$ & $0.121$ \\
$400\times50$ & $34$ & $0.021$ & $0.0037$ & $0.0026$ & $0.008$ & $0.0007$ & $0.0005$ & $160.2$ & $160.8$ & $159.5$ & $0.099$ & $0.097$ \\
$480\times60$ & $30$ & $0.074$ & $0.0075$ & $0.0052$ & $0.041$ & $0.0013$ & $0.0009$ & $160.9$ & $160.9$ & $160.3$ & $0.100$ & $0.091$ \\
$560\times70$ & $39$ & $0.055$ & $0.0055$ & $0.0038$ & $0.017$ & $0.0010$ & $0.0007$ & $161.1$ & $161.5$ & $160.7$ & $0.081$ & $0.077$ \\
$640\times80$ & $39$ & $0.033$ & $0.0041$ & $0.0028$ & $0.052$ & $0.0013$ & $0.0009$ & $161.4$ & $162.3$ & $160.9$ & $0.072$ & $0.070$ \\
$720\times90$ & $39$ & $0.041$ & $0.0050$ & $0.0035$ & $0.035$ & $0.0010$ & $0.0007$ & $161.9$ & $163.3$ & $161.1$ & $0.065$ & $0.063$ \\
$800\times100$ & $44$ & $0.073$ & $0.0050$ & $0.0035$ & $0.041$ & $0.0011$ & $0.0008$ & $162.0$ & $163.4$ & $161.4$ & $0.058$ & $0.055$ \\
\midrule
\multicolumn{13}{@{}l}{\textit{Yuksel (stage 2)}} \\
$160\times20$ & $179$ & $0.030$ & $0.0033$ & $0.0023$ & $0.001$ & $0.0001$ & $0.0001$ & $157.7$ & $158.3$ & $156.9$ & $0.077$ & $0.080$ \\
$240\times30$ & $201$ & $0.100$ & $0.0048$ & $0.0034$ & $0.001$ & $0.0000$ & $0.0000$ & $160.0$ & $160.1$ & $159.3$ & $0.036$ & $0.046$ \\
$320\times40$ & $394$ & $0.091$ & $0.0040$ & $0.0028$ & $0.021$ & $0.0007$ & $0.0005$ & $160.9$ & $160.9$ & $160.7$ & $0.026$ & $0.035$ \\
$400\times50$ & $414$ & $0.086$ & $0.0027$ & $0.0019$ & $0.013$ & $0.0006$ & $0.0004$ & $159.2$ & $160.0$ & $160.0$ & $0.029$ & $0.037$ \\
$480\times60$ & $884$ & $0.031$ & $0.0009$ & $0.0006$ & $0.016$ & $0.0008$ & $0.0005$ & $159.7$ & $160.5$ & $160.5$ & $0.024$ & $0.029$ \\
$560\times70$ & $1119$ & $0.100$ & $0.0019$ & $0.0014$ & $0.049$ & $0.0012$ & $0.0009$ & $159.5$ & $160.3$ & $160.3$ & $0.019$ & $0.023$ \\
$640\times80$ & $1749$ & $0.100$ & $0.0019$ & $0.0013$ & $0.078$ & $0.0012$ & $0.0009$ & $158.9$ & $159.7$ & $159.7$ & $0.012$ & $0.014$ \\
$720\times90$ & $1374$ & $0.100$ & $0.0018$ & $0.0013$ & $0.100$ & $0.0020$ & $0.0014$ & $159.5$ & $160.3$ & $160.3$ & $0.012$ & $0.015$ \\
$800\times100$ & $1329$ & $0.100$ & $0.0018$ & $0.0012$ & $0.100$ & $0.0014$ & $0.0010$ & $159.5$ & $160.3$ & $160.3$ & $0.009$ & $0.013$ \\
\midrule
\multicolumn{13}{@{}l}{\textit{Du--Olhoff}} \\
$160\times20$ & $53$ & $0.056$ & $0.0055$ & $0.0038$ & $0.013$ & $0.0007$ & $0.0005$ & $168.6$ & $169.2$ & $169.2$ & $0.126$ & $0.116$ \\
$240\times30$ & $34$ & $0.073$ & $0.0064$ & $0.0044$ & $0.007$ & $0.0009$ & $0.0006$ & $167.5$ & $167.7$ & $167.3$ & $0.127$ & $0.123$ \\
$320\times40$ & $59$ & $0.049$ & $0.0053$ & $0.0036$ & $0.006$ & $0.0008$ & $0.0005$ & $166.3$ & $166.3$ & $165.7$ & $0.151$ & $0.142$ \\
$400\times50$ & $52$ & $0.028$ & $0.0043$ & $0.0030$ & $0.003$ & $0.0005$ & $0.0003$ & $167.0$ & $167.0$ & $166.3$ & $0.132$ & $0.123$ \\
$800\times100$ & $232$ & $0.011$ & $0.0015$ & $0.0010$ & $0.004$ & $0.0006$ & $0.0004$ & $165.4$ & $165.4$ & $165.4$ & $0.169$ & $0.159$ \\
\bottomrule
\end{tabular}
\end{table}

\begin{table}[p]
\centering\small
\caption{Candidate common rules replayed on every recorded history. Entries: stop
  iteration ($\omega_1$ loss vs end state, \%). Yuksel: rule applied to stage~2
  only (stage-1 handoff kept native). Objective and $\Mnd$ ranges are taken over
  a window of $W$ iterations as in Eq.~\eqref{eq:R1}.}
\label{tab:C}
\begin{tabular}{@{}lrcccccc@{}}
\toprule
mesh & $\kstag$
  & \makecell{\textbf{R1}: obj.\ $< 10^{-3}$\\and $\Mnd < 5\cdot10^{-3}$, $W{=}10$}
  & \makecell{R1 tight: $5\cdot10^{-4}$\\/ $5\cdot10^{-3}$, $W{=}10$}
  & \makecell{obj.\ $< 10^{-3}$\\only, $W{=}10$}
  & \makecell{R2: $\maxdx \le 0.04$\\for 10 consecutive it.}
  & \makecell{$\maxdx \le 0.025$\\(single it.)}
  & \makecell{$\Mnd < 5\cdot10^{-3}$\\only, $W{=}10$} \\
\midrule
\multicolumn{8}{@{}l}{\textit{Proposed}} \\
$160\times20$ & $29$ & 44 $(-0.24)$ & 48 $(-0.24)$ & 44 $(-0.24)$ & 41 $(-0.22)$ & 36 $(-0.21)$ & 41 $(-0.22)$ \\
$240\times30$ & $95$ & 46 $(+0.33)$ & 48 $(+0.33)$ & 46 $(+0.33)$ & 45 $(+0.33)$ & 42 $(+0.33)$ & 39 $(+0.32)$ \\
$320\times40$ & $30$ & 54 $(-0.13)$ & 56 $(-0.13)$ & 54 $(-0.13)$ & 44 $(-0.14)$ & 41 $(-0.16)$ & 41 $(-0.16)$ \\
$400\times50$ & $34$ & 91 $(-0.09)$ & 132 $(-0.04)$ & 91 $(-0.09)$ & 41 $(-0.32)$ & 37 $(-0.42)$ & 36 $(-0.42)$ \\
$480\times60$ & $30$ & 62 $(+0.03)$ & 71 $(+0.03)$ & 62 $(+0.03)$ & 59 $(+0.02)$ & 59 $(+0.02)$ & 42 $(-0.09)$ \\
$560\times70$ & $39$ & 64 $(-0.05)$ & 68 $(-0.04)$ & 64 $(-0.05)$ & 76 $(-0.05)$ & 69 $(-0.04)$ & 42 $(-0.26)$ \\
$640\times80$ & $39$ & 67 $(-0.03)$ & 72 $(-0.03)$ & 67 $(-0.03)$ & 56 $(-0.06)$ & 55 $(-0.06)$ & 38 $(-0.57)$ \\
$720\times90$ & $39$ & 82 $(+0.02)$ & 96 $(+0.04)$ & 82 $(+0.02)$ & 68 $(-0.03)$ & 60 $(-0.06)$ & 36 $(-0.76)$ \\
$800\times100$ & $44$ & 98 $(+0.08)$ & 106 $(+0.09)$ & 98 $(+0.08)$ & \textit{none} & 126 $(+0.09)$ & 42 $(-0.56)$ \\
\midrule
\multicolumn{8}{@{}l}{\textit{Yuksel (stage 2)}} \\
$160\times20$ & $179$ & 221 $(-0.22)$ & 225 $(-0.21)$ & 221 $(-0.22)$ & 176 $(-0.53)$ & 208 $(-0.26)$ & 147 $(-0.88)$ \\
$240\times30$ & $201$ & 208 $(-0.46)$ & 263 $(-0.10)$ & 208 $(-0.46)$ & 229 $(-0.39)$ & 223 $(-0.40)$ & 193 $(-0.26)$ \\
$320\times40$ & $394$ & 364 $(-0.15)$ & 364 $(-0.15)$ & 364 $(-0.15)$ & 449 $(+0.13)$ & 473 $(+0.11)$ & 273 $(+2.36)$ \\
$400\times50$ & $414$ & 447 $(+0.33)$ & 533 $(+0.15)$ & 447 $(+0.33)$ & 409 $(+0.53)$ & 671 $(+0.06)$ & 335 $(+2.11)$ \\
$480\times60$ & $884$ & 718 $(+0.91)$ & 788 $(+0.72)$ & 718 $(+0.91)$ & 790 $(+0.71)$ & 784 $(+0.72)$ & 607 $(+2.62)$ \\
$560\times70$ & $1119$ & 1010 $(+0.81)$ & 1235 $(+0.23)$ & 1010 $(+0.81)$ & 1291 $(+0.13)$ & 1191 $(+0.30)$ & 853 $(+2.71)$ \\
$640\times80$ & $1749$ & 1719 $(+0.61)$ & 1770 $(+0.45)$ & 1719 $(+0.61)$ & \textit{none} & \textit{none} & 1563 $(+2.40)$ \\
$720\times90$ & $1374$ & 1325 $(+0.65)$ & 1389 $(+0.47)$ & 1325 $(+0.65)$ & \textit{none} & \textit{none} & 1182 $(+2.42)$ \\
$800\times100$ & $1329$ & 1292 $(+0.64)$ & 1342 $(+0.47)$ & 1292 $(+0.64)$ & \textit{none} & \textit{none} & 1171 $(+2.28)$ \\
\midrule
\multicolumn{8}{@{}l}{\textit{Du--Olhoff}} \\
$160\times20$ & $53$ & 66 $(+0.27)$ & 77 $(+0.22)$ & 66 $(+0.27)$ & 81 $(+0.23)$ & 72 $(+0.21)$ & 62 $(+0.29)$ \\
$240\times30$ & $34$ & 57 $(-0.05)$ & 60 $(-0.05)$ & 46 $(-0.16)$ & 70 $(-0.04)$ & 64 $(-0.05)$ & 47 $(-0.12)$ \\
$320\times40$ & $59$ & 78 $(-0.10)$ & 81 $(-0.10)$ & 65 $(-0.28)$ & 90 $(-0.10)$ & 69 $(-0.18)$ & 71 $(-0.15)$ \\
$400\times50$ & $52$ & 78 $(-0.13)$ & 84 $(-0.11)$ & 59 $(-0.35)$ & 69 $(-0.21)$ & 64 $(-0.30)$ & 67 $(-0.25)$ \\
$800\times100$ & $232$ & 241 $(-0.04)$ & 241 $(-0.04)$ & 122 $(+1.09)$ & 105 $(+1.31)$ & 130 $(+1.05)$ & 241 $(-0.04)$ \\
\bottomrule
\end{tabular}
\end{table}
\end{landscape}

\begin{table}[htbp]
\centering\small
\caption{Summary over all 23 recorded histories (methods $\times$ meshes),
  sorted by worst $\omega_1$ loss. Worst loss = largest $\omega_1$ deficit at the
  stop relative to the end state; ratio = $k^*/\kstag$ (1 = stops exactly at
  stagnation, $<1$ = before, $>1$ = after). ``obj.'' and ``$\Mnd$'' are window
  ranges as in Eq.~\eqref{eq:R1}; rules without $W$ act on a single iteration.
  ``flip'' = fraction of elements crossing 0.5.}
\label{tab:D}
\begin{tabular}{@{}lcccc@{}}
\toprule
rule & \makecell{never\\fired} & \makecell{worst $\omega_1$\\loss [\%]}
  & \makecell{worst $\Mnd$\\excess} & \makecell{$k^*/\kstag$\\min / median / max} \\
\midrule
$\maxdx \le 0.01$                                       & 7 & 0.03 & $+0.005$ & 1.06 / 2.39 / 7.62 \\
obj.\ $<5\cdot10^{-4}$ \& $\Mnd<5\cdot10^{-3}$, $W{=}20$ & 0 & 0.33 & $+0.011$ & 0.71 / 1.58 / 5.79 \\
obj.\ $<5\cdot10^{-4}$ \& $\Mnd<5\cdot10^{-3}$, $W{=}10$ & 0 & 0.72 & $+0.010$ & 0.51 / 1.37 / 3.88 \\
$\relL < 10^{-3}$                                       & 2 & 0.72 & $+0.003$ & 0.88 / 2.02 / 9.23 \\
$\RMS(\dx) < 8.8\cdot10^{-4}$ (Du--Olhoff $c{=}0.05$)    & 1 & 0.72 & $+0.005$ & 0.88 / 1.75 / 8.57 \\
\textbf{R1}: obj.\ $<10^{-3}$ \& $\Mnd<5\cdot10^{-3}$, $W{=}10$ & 0 & 0.91 & $+0.010$ & 0.48 / 1.25 / 2.68 \\
$\maxdx \le 0.02$, $W{=}10$                              & 6 & 0.98 & $+0.094$ & 0.58 / 1.71 / 3.53 \\
obj.\ $<10^{-3}$ only, $W{=}20$                          & 0 & 1.01 & $+0.104$ & 0.60 / 1.49 / 4.15 \\
obj.\ $<5\cdot10^{-4}$ only, $W{=}10$                    & 0 & 1.03 & $+0.108$ & 0.51 / 1.37 / 3.88 \\
$\maxdx \le 0.025$                                      & 3 & 1.05 & $+0.111$ & 0.44 / 1.24 / 2.86 \\
obj.\ $<10^{-3}$ only, $W{=}10$                          & 0 & 1.09 & $+0.118$ & 0.48 / 1.13 / 2.68 \\
obj.\ $<10^{-3}$ \& $\Mnd<10^{-2}$, $W{=}10$             & 0 & 1.09 & $+0.118$ & 0.48 / 1.25 / 2.68 \\
$\relL < 2\cdot10^{-3}$                                 & 0 & 1.22 & $+0.009$ & 0.84 / 1.42 / 8.07 \\
R2: $\maxdx \le 0.04$, $W{=}10$                          & 4 & 1.31 & $+0.134$ & 0.45 / 1.33 / 2.06 \\
$\omega_1$ rel.\ range $<10^{-3}$, $W{=}10$              & 0 & 1.51 & $+0.118$ & 0.43 / 1.10 / 1.67 \\
$\Mnd < 5\cdot10^{-3}$ only, $W{=}10$                    & 0 & 2.71 & $+0.008$ & 0.41 / 0.96 / 1.41 \\
flip $<10^{-3}$, $W{=}5$                                 & 0 & 2.82 & $+0.084$ & 0.60 / 1.27 / 2.93 \\
flip $<10^{-3}$                                         & 0 & 3.29 & $+0.102$ & 0.36 / 0.94 / 1.57 \\
$\relL < 5\cdot10^{-3}$                                 & 0 & 7.35 & $+0.328$ & 0.16 / 1.09 / 1.70 \\
$\RMS(\dx) < 3.5\cdot10^{-3}$ ($c{=}0.2$)                & 0 & 7.35 & $+0.328$ & 0.16 / 1.09 / 1.67 \\
$\maxdx \le 0.04$                                       & 0 & 9.27 & $+0.405$ & 0.08 / 1.14 / 1.79 \\
$\relL < 10^{-2}$                                       & 0 & 9.67 & $+0.418$ & 0.07 / 0.91 / 1.27 \\
\bottomrule
\end{tabular}
\end{table}

%======================================================================
\section{Files}

\begingroup\raggedright
\begin{itemize}
\item \file{scripts/\allowbreak{}stopstudy\_\allowbreak{}\{proposed,yuksel,olhoff\}.m} --- recording
  runners; they only set existing configuration fields and run the production
  code paths.
\item \file{scripts/\allowbreak{}stopstudy\_\allowbreak{}metrics.m}, \file{scripts/\allowbreak{}replay.py} (rule replay
  and $\kstag$), \file{scripts/\allowbreak{}make\_\allowbreak{}tables.py}, \file{scripts/\allowbreak{}plots.py},
  \file{scripts/\allowbreak{}plot\_\allowbreak{}summary.py}.
\item \file{data/\allowbreak{}<method>\_\allowbreak{}<mesh>.csv} (per-iteration metrics),
  \file{data/\allowbreak{}replay\_\allowbreak{}all.csv}, \file{data/\allowbreak{}tables.md}.
\item \file{fig/\allowbreak{}traces\_\allowbreak{}\{proposed,yuksel,olhoff\}.png}
  (Figs.~\ref{fig:traces-proposed}--\ref{fig:traces-olhoff}),
  \file{fig/\allowbreak{}stop\_\allowbreak{}rules\_\allowbreak{}summary.png} (Fig.~\ref{fig:summary}).
\end{itemize}
\endgroup

\end{document}
